\input format.tex

%\usepackage{arydshln}
\usepackage{graphicx}
\graphicspath{{cores/}}

\begin{document}

%%%
\setlength{\arrayrulewidth}{0.5pt}
\fontsize{9.3pt}{17pt}\selectfont
\color{gray2}

\vspace*{0mm}
\begin{center}
%{\bf\sanhao 肠道致病菌}
{\bf\sanhao 致病菌}
\end{center}

\medskip

\noindent
{\bf\xiaosihao Tom先生：}


\bigskip

%您的肠道内未检测到可能导致疾病的致病菌，这是个好消息。当您的免疫力较强，可持续抵御致病菌的入侵。但请注意，当您的免疫力较低，肠道致病菌可能会增加，早期可能引起腹泻、腹痛、呕吐、感染等病症。
%您的肠道中检测出可能导致疾病的致病菌2种，含量尚处于正常范围。但请注意，当免疫功能低下时，它们可能引起相关疾病。建议您持续监测。下面的检测列表可以帮助您快速锁定可能的原因。
%您的肠道中检测出可能导致疾病的致病菌2种，其中病原性大肠埃希氏菌、脆弱拟杆菌含量偏高。致病菌可能引起腹泻、腹痛、呕吐、肠炎等相关疾病。建议您持续监测，并到正规医院咨询专科医生，遵从医生建议做进一步检查。

您的肠道内
检测出2种致病菌，其中病原性大肠埃希氏菌、脆弱拟杆菌含量高于参考人群，当身体免疫力较强时可能不会出现不适感。但请注意，当您的免疫力较弱时，初期可能引起腹泻、腹痛、感染等病症，长期可能导致上表中所述疾病。建议您持续监测，必要时请到正规医院咨询专科医生，遵从医生的建议做进一步检查。

\fontsize{9.3pt}{9.8pt}\selectfont
\bigskip
{
\lantxh
%\begin{tctabularx}{tabularx={m{3.75cm}m{4.9cm}<{\centering}m{1.22cm}<{\centering}m{1.8cm}<{\centering}m{2cm}<{\centering}}}
\begin{tctabularx}{tabularx={m{4.8cm}m{5.2cm}<{\centering}m{4.55cm}<{\centering}}}
%\cellcolor{kelly} &
%\cellcolor{kelly} &
%\cellcolor{kelly} &
%\cellcolor{kelly} &
%\cellcolor{kelly}
&&
\\[-6pt]
\cellcolor{kelly} \raisebox{2.614pt}{\color{black70}\bf 测试项目} &
\cellcolor{kelly} \raisebox{2.614pt}{\color{black70}\bf 分布} &
%\cellcolor{kelly} \raisebox{2.614pt}{\color{black70}\bf 检测值} &
%\cellcolor{kelly} \raisebox{2.614pt}{\color{black70}\bf 参考范围} &
\cellcolor{kelly} \raisebox{2.614pt}{\color{black70}\bf 检测结果}
%{\color{black70}\bf 测试项目} & {\color{black70}\bf 分布} & {\color{black70}\bf 检测值} &{\color{black70}\bf 参考范围} & {\color{black70}\bf 检测结果}
\end{tctabularx}

\vspace*{-4.25mm}
\fontsize{8.8pt}{11pt}\selectfont
%\begin{longtable}{m{3.75cm}m{5cm}<{\centering}m{0cm}@{}m{1.22cm}<{\centering}m{1.8cm}<{\centering}m{2cm}<{\centering}}
\begin{longtable}{m{4.8cm}m{5.2cm}<{\centering}m{0cm}@{}m{4.61cm}<{\centering}}
\hline
\parbox[c]{\hsize}{\vskip6pt 病原性大肠埃希氏菌\\可能导致肠胃炎、尿路感染、新生儿脑膜炎、腹膜炎等 \vskip6pt} & \parbox[c]{\hsize}{\vskip6pt\centerline{\raisebox{-1.5ex}{\includegraphics[width=5cm,height=0.3cm]{zxdachangganjun_dis.png}}}\vskip6pt}  &\hspace*{-4.273390214cm}\raisebox{-0.45ex}{\includegraphics{triangle.pdf}} & \begin{minipage}{4.60cm}\begin{center}{高\\ \bahao 可能增加肠胃炎、尿路感染、新生儿脑膜炎、腹膜炎等的风险 }\end{center} \end{minipage} \\
\hline
\parbox[c]{\hsize}{\vskip6pt 脆弱拟杆菌\\可能导致菌血症、腹内感染、腹膜炎 \vskip6pt} & \parbox[c]{\hsize}{\vskip6pt\centerline{\raisebox{-1.5ex}{\includegraphics[width=5cm,height=0.3cm]{zxcuiruoniganjun_dis.png}}}\vskip6pt}  &\hspace*{-3.783564884cm}\raisebox{-0.45ex}{\includegraphics{triangle.pdf}} & \begin{minipage}{4.60cm}\begin{center}{高\\ \bahao 可能增加菌血症、腹内感染、腹膜炎的风险 }\end{center} \end{minipage} \\
\hline
\end{longtable}

\noindent
\tikz\draw[green2,fill=green2](0,0) circle[radius=4pt]; 区域表示正常人群的一般水平 \tikz\draw[darkblue,fill=darkblue](0,0) circle[radius=4pt]; 区域表示仍在正常水平，但较一般水平偏高或偏低 %\tikz\draw[red,fill=red](0,0) circle[radius=4pt]; 而红色表示为有风险
}

\bigskip
{\centering {\includegraphics{jieshi.pdf}}}
\end{document}
